This appendix collects background for a future formal theory of arithmetic
expressions, continuing Subsection~\ref{subsec:p0-evaluation-effectivity}.
The constructions of Bauer--Taylor and Conway are cited external mathematics;
their proposed use in AEG is a research direction.  No general convergence
theory, new expression calculus, or additional mechanized theorem is asserted
here.

\subsection{Bauer--Taylor: cuts, observable predicates, and interval information}
\label{subsec:p0-taylor-cuts}

Bauer and Taylor construct the real line in Abstract Stone Duality (ASD)
using two-sided Dedekind cuts \cite[Sections 2 and 6]{BauerTaylor2009Dedekind}.
In the familiar rational presentation, a real $x$ has lower and upper predicates
\[
  \delta(q)\equiv q<x,\qquad \upsilon(q)\equiv x<q,
  \qquad q\in\Q.
\]
Conversely, inhabited, rounded lower and upper predicates that are disjoint
and satisfy locatedness specify a real.  Roundedness means
\[
 \delta(q)\iff\exists r\in\Q\;(q<r\ \land\ \delta(r)),\qquad
 \upsilon(q)\iff\exists r\in\Q\;(\upsilon(r)\ \land\ r<q),
\]
and locatedness is the overlapping alternative
\[
  p<q\quad\Longrightarrow\quad\delta(p)\lor\upsilon(q).
\]
This is not a Boolean decision of $q<x$ versus $x<q$ at the same rational $q$.
In particular, both strict predicates fail when $q=x$ is rational.

Taylor's computational reading uses predicates valued in the Sierpi\'nski
space $\Sigma$, rather than Boolean tests: a positive observation may be
witnessed by termination, while absence of a witness does not decide its
negation \cite[Section 2]{Taylor2006IntervalWithoutIntervals}.
Dropping locatedness allows the two sides to leave an interval of candidates.
Additional lower and upper evidence narrows that interval: inclusion of
evidence corresponds to reverse inclusion of candidate intervals.  This
\emph{information order} differs from numerical order.

Taylor develops a calculus for exact real analysis and translates it to
interval computation \cite[Sections 5--6]{Taylor2006IntervalWithoutIntervals}.
Thus an exact denotation and its finite enclosure need not be the same object.
The order-theoretic construction also permits other countable dense bases,
including dyadics and real algebraic numbers; an order isomorphism alone does
not preserve arithmetic.  The interpolation discussion explicitly mentions
Conway's simplest interpolant as an alternative to an arithmetic midpoint
\cite[Section 6]{BauerTaylor2009Dedekind}.

\subsection{Conway: recursive options and simplicity}
\label{subsec:p0-conway-cuts}

Conway forms a number $\{L\mid R\}$ from sets of previously constructed
numbers with $\ell<r$ for every $\ell\in L$, $r\in R$.  Either set may be
empty.  The simplicity theorem identifies its value as the unique simplest
surreal satisfying those strict inequalities, where simplicity belongs to
the recursive construction, not to numerical magnitude
\cite[Chapters 0--2]{Conway1976NumbersGames}.  For example,
\[
  \{\mid\}=0,\qquad \{0\mid\}=1,\qquad
  \{0\mid1\}=\tfrac12,\qquad \{0\mid3\}=1.
\]
The last form selects $1$; it does not retain an unknown target somewhere
between $0$ and $3$, and does not select the arithmetic midpoint $3/2$.

Arithmetic is itself recursively defined.  For example, using $x^L,x^R$
and $y^L,y^R$ for typical options,
\[
 x+y=\{x^L+y,\ x+y^L\mid x^R+y,\ x+y^R\}.
\]
All available options of each indicated kind are included.  Compatibility
with equality and the field laws require the accompanying proofs; they do
not follow from bracket notation alone.  Transfinite construction yields
the proper class $\mathbf{No}$, a real-closed ordered field containing the
reals, infinite numbers, and nonzero infinitesimals.  For instance,
\[
 \omega=\{0,1,2,\ldots\mid\},\qquad
 \omega^{-1}=\{0\mid1,\tfrac12,\tfrac14,\ldots\}.
\]
An infinitesimal here is a definite nonzero number, not an unresolved
approximation error \cite{Ehrlich2011ConwayNames}.

\subsection{Finite names, constructions, and observations}
\label{subsec:p0-cut-representation-boundaries}

Finite well-founded option trees in Conway's primitive cut language yield
exactly the dyadic rationals $\mathbb Z[1/2]$.  Nevertheless, in the expanded
field language $1/3$ has a finite expression, although its primitive cut
construction is infinite \cite[Chapter 2]{Conway1976NumbersGames}.
Finite syntax, finite unfolding, and terminating computation are therefore
separate requirements.  A finite rule for an infinite cut must state its
semantics and any effective approximation guarantees.  Over a fixed
countable alphabet, finite names cannot cover every real without additional
parameters; neither construction removes this limitation.

The ambient carrier also matters.  In $\R$, all rational lower and upper
comparisons determine a unique point.  In $\mathbf{No}$, an embedded
\emph{irrational} real $r$ and $r+\omega^{-1}$ have identical strict
comparisons with every rational:
the absolute value of each such real difference exceeds the positive
infinitesimal.  Rational
comparisons alone consequently do not distinguish them; Conway's simplest
choice supplies additional structure.  Moreover, cuts require an ordered
setting, whereas the ordinary expression syntax over a general field does
not.  They cannot be silently imposed on complex or positive-characteristic
fields.

\subsection{The Adva through interface and the next formal questions}
\label{subsec:p0-cut-formal-programme}

Adva's finite research adapter gives the middle bar a different, explicitly
typed interpretation \cite{Adva2026GroundedThrough}.  For maps to a declared
middle object $N$, the through relation is the fibre product
\[
 \langle L\mid_N R\rangle=L\times_N R
   =\{(\ell,r):q_L(\ell)=q_R(r)\}.
\]
In the implemented adapter the specializations may instead be relations;
a pair is retained when the two middle-image sets intersect.  The middle
object is a declared finite quotient by exact upper-cut wire indices,
derived from checked program incidences.  Copy may leave several matches;
the complete program slice and occurrence distinctions remain attached.
This is bounded research evidence, not a stable native cut operation, a
Conway simplicity selector, or an implementation of ASD.

There is already a useful obstruction to identifying these interpretations
\cite{Adva2026CutObjectification}.  Write $\operatorname{Obj}$ for Conway
value selection and let $g(z)=3z$, acting on each presented option.  Then
\[
 g\bigl(\operatorname{Obj}\{0\mid1\}\bigr)=\tfrac32,
 \qquad
 \operatorname{Obj}\{0\mid3\}=1.
\]
Scaling the options is not the recursively defined multiplication of the
number.  A value projection therefore cannot be assumed to commute with
arbitrary transformations of presentations.

For AEG the proposed next step is to specify, separately:
\begin{enumerate}
 \item a finite expression signature, constant names, and partial ordinary
       interpretation with every division guard retained;
 \item constructors for exact denotations, with explicit root, branch,
       embedding, or cut conditions when required;
 \item observations that produce justified bounds while retaining shared
       dependencies and any symbolically unexpanded remainder;
 \item distinct judgments for presentation equality, denotational equality,
       information refinement, and typed through compatibility.
\end{enumerate}
These are open design and proof obligations.  The comparison proposes no
identification of Conway simplicity, AEG history, and computational resource
cost, and no extension of the formal certificates already recorded for this
manuscript.
